
\section{Introdução}
\label{sec:intro}

O desenvolvimento de Veículos Aéreos Não Tripulados (VANTs) ou do termo em inglês, Unmanned Aerial Vehicles (UAVs), popularmente conhecidos como \textit{drones}, tem crescido exponencialmente nas últimas décadas, impulsionado pelo avanço em áreas como eletrônica embarcada, sensores, comunicação sem fio e algoritmos de controle. Inicialmente restritos a aplicações militares, os VANTs tornaram-se elementos fundamentais em um vasto conjunto de contextos civis e científicos, como monitoramento ambiental, mapeamento geográfico, inspeções industriais, agricultura de precisão e resposta a desastres \citep{chung2018survey, hong2020swarm}. Essa expansão foi acompanhada por um aumento da autonomia e da capacidade de processamento dos sistemas embarcados, permitindo o surgimento de arquiteturas de controle mais sofisticadas e cooperativas.

Entretanto, apesar dos avanços, a operação isolada de um único VANT apresenta limitações inerentes, especialmente em termos de cobertura espacial, redundância e tolerância a falhas. Essas restrições motivaram a transição para abordagens cooperativas envolvendo múltiplos veículos, conhecidas como \textit{enxames de VANTs} (\textit{UAV swarms}) \citep{bu2024advancement}. Nesses sistemas, diversos VANTs atuam de maneira coordenada e distribuída, compartilhando informações e ajustando seus comportamentos locais de modo a atingir objetivos coletivos de forma mais eficiente e robusta do que um agente individual seria capaz de realizar.

O princípio fundamental dos enxames de VANTs inspira-se na \textit{inteligência coletiva} observada em sistemas naturais, como colônias de insetos, bandos de aves e cardumes de peixes. Essa abordagem, conhecida como \textit{Inteligência de Enxame} (\textit{Swarm Intelligence}), baseia-se em interações locais simples entre agentes, que, quando combinadas, dão origem a comportamentos globais emergentes \citep{ren2007consensus, chung2018survey}. Em um contexto de veículos aéreos, isso se traduz na capacidade de auto-organização, reconfiguração dinâmica e coordenação cooperativa, mesmo sob condições de comunicação limitada ou ambientes parcialmente desconhecidos.

Nesse cenário, o projeto de arquiteturas de controle adequadas é um dos principais desafios de pesquisa. Tais arquiteturas devem integrar comunicação eficiente, mecanismos de consenso, estratégias de controle de formação e coordenação de missão, equilibrando robustez, escalabilidade e custo computacional \citep{ouyang2023formation, bu2024advancement}. Além disso, a crescente demanda por sistemas multiagente em aplicações reais — como vigilância cooperativa, entrega autônoma e busca e salvamento — reforça a necessidade de modelos de simulação e validação que permitam avaliar o comportamento coletivo antes da implementação em campo.

Dessa forma, este trabalho tem como objetivo apresentar uma análise abrangente das principais abordagens, algoritmos e arquiteturas aplicadas ao controle cooperativo de enxames de VANTs. São discutidos conceitos fundamentais, como os níveis de autonomia e os princípios de sistemas multiagente, bem como metodologias de formação, estratégias de comunicação e desafios relacionados à coordenação distribuída. 

Com base nesses fundamentos, propõe-se investigar, por meio de simulações e testes comparativos, o desempenho de diferentes algoritmos de controle de formação — como \textit{Leader–Follower}, \textit{Consensus-Based}, \textit{Behavior-Based}, \textit{Virtual Structure} e \textit{Artificial Potential Fields (APF)} \citep{bu2024advancement} — em diferentes cenários de operação. O estudo busca avaliar criticamente a aplicabilidade e a eficiência de cada método sob condições diversas, incluindo ambientes com obstáculos, variações de topologia de comunicação e ajustes dinâmicos de parâmetros.

Por fim, será investigado o potencial de métodos de hibridização entre algoritmos de controle, buscando combinar características complementares de cada abordagem. A intenção é identificar configurações capazes de aprimorar a estabilidade da formação e ampliar a adaptabilidade do enxame frente a distintos contextos operacionais. Esse esforço visa consolidar uma arquitetura de sistema que sustente a coordenação autônoma, eficiente e robusta de múltiplos drones em cenários variados.


